% V. Методика на експерименталната проверка.
% Този раздел заявява ПЛАНА на експеримента преди да има измервания.
% В него няма и не трябва да има числа.
\section{Методика на експерименталната проверка}
\label{sec:method}

Измерването отговаря на два отделни въпроса и затова е разделено на две части. Първата
е за коректност: дава ли всяка стъпка същия резултат като еталонната реализация на
същия алгоритъм. Втората е за цена: колко време и колко памет отнема стъпката. Двете не
се смесват, защото по-бърза стъпка с различен резултат не е по-бърза, а друга стъпка.

Еталон за сравнение са реализациите на същите стъпки в утвърдените библиотеки с отворен
код~\cite{rusu2011pcl,bradski2000opencv}. За заявката за най-близки съседи еталонът е
изчерпателното търсене, чийто резултат е точен по определение и служи за проверка на
индекса. Съвпадението се проверява по множеството от върнати индекси, а при равни
разстояния се допуска различен ред, защото редът не е част от договора.

Измерваните величини са: време за построяване на индекса, време за заявка при няколко
стойности на $k$ и на радиуса, изразходвана памет от индекса, и време на всяка стъпка
от конвейера поотделно. Всяко измерване се повтаря многократно и се отчита медианата
заедно с разсейването, защото средното е чувствително към единични отклонения от
планировчика на операционната система.

\TODO{набор от данни: източник, брой облаци, брой точки, начин на снемане. Използват
се само публично достъпни набори и не се съобщава характеристика, която не е проверена
в оригиналния източник}

\TODO{конфигурация на машината: процесор, набор инструкции, размер на кеша, памет,
компилатор и версия, флагове за оптимизация}

\TODO{протокол на изпълнение: брой повторения, загряващи изпълнения, състояние на
мащабирането на честотата}

Таблиците с резултати имат предварително фиксирана форма, за да не бъде изборът им
подчинен на това, което измерването е показало. Таблица~1 съпоставя двете реализации на
индекса по време за построяване, време за заявка и памет. Таблица~2 дава времето на
всяка стъпка от конвейера при скаларна и при векторизирана реализация. Таблица~3
съпоставя резултата на всяка стъпка с еталона по избраната за нея мярка за съвпадение.
Всички клетки остават празни до провеждането на измерването.
